%======================================================================
\section{Causal services and gadget channels}\label{sec:channel}\label{sec:channel-setting}
%======================================================================

Part~II can be read on its own by a reader who comes from quantum information. From Part~I it needs three definitions, all in
Section~\ref{sec:criterion}. A \emph{chunk} of a group is a finite subset containing $1$, with the partial multiplication it inherits. A
\emph{sofic model} of the chunk at accuracy $1/r$ assigns a permutation of a finite set to each element, with products correct and distinct
elements separated on all but a fraction $1/r$ of the points. The \emph{profile} $\sigma_E(r)$ is the least number of points of such a model.
The configuration lamps enter mainly through two facts, quoted where they are used: their pair areas (Theorem~\ref{thm:pair-area}) and their
profile, Theorem~A in the precise form of Theorem~\ref{thm:main-group}. Concretely, the lower bounds use one black box, the transport of
Section~\ref{sec:lamp-channel}: for each $N$ there are $2^N$ words of length at most $63N^2$ whose copies of $S_3$ commute pairwise at area at
most $C_1N^\kappa$, with $\kappa=6+\log107$ (the introduction says where $107$ comes from). The upper bounds use another, the hash models
behind Theorem~A(b), which give sofic models with $\exp(O(r\log r))$ points.

A finite set of relations of a group becomes a quantum channel of fixed dimension. Each relation contributes a few small unitary blocks; the
channel evaluates the blocks in the regular representation of the group and traces out the group von Neumann algebra. Every channel on
$\mathbb C^d$ has a Stinespring dilation, a unitary on the input and a pure environment of dimension at most $d^2$ that is discarded after
use. Applied $n$ times with a fresh environment each time, the channel needs almost no memory, but the discarded environments carry away
entropy that fresh pure qubits paid for. We charge for that purity, and the memory is never reset. The question is how much memory a device
needs when its purity is far below what the Stinespring device spends.

In~\cite{DouglasMghirbiB} we showed that the memory needed to apply such a channel repeatedly, releasing each output before the next input
arrives, is bounded below by the size of approximate representations of its relations. Part~II sharpens that link, so that it loses a single
square root, and applies it to the configuration lamps. The outcome is Theorem~B, stated precisely as Theorem~\ref{thm:main-channel}: an
explicit channel in the closure of finite tracial channels whose repeated use with logarithmic purity needs memory $n^{1-o(1)}$, while memory
$O(n\log n)$ suffices. The lower bounds come from the sharpened link, the weighted one-round theorem of Section~\ref{sec:weighted-statement}.
The upper bounds come from the sofic models of $\Gamma$, which give devices directly (Proposition~\ref{prop:sofic-upper}) and, at the rank
rate, through Theorem~\ref{thm:streaming}.

We recall the framework of~\cite{DouglasMghirbiB} and keep its notation.

\paragraph{Channels.} Channels act on $M_d=M_d(\mathbb C)$. The normalized Choi state of a channel $\Psi$ is
$\mathcal J(\Psi)=(\Psi\otimes\mathrm{id})(|\Omega\rangle\langle\Omega|)$, with $|\Omega\rangle=d^{-1/2}\sum_i|ii\rangle$, and
$\ddia(\Psi,\Psi')=\frac12\|\Psi-\Psi'\|_\diamond$ is the diamond distance. We write $\tau_M$ for the normalized trace on $M_M$. A channel $\Psi$ has a
\emph{finite tracial factorization} with bath dimension $M$ if $\Psi(X)=(\mathrm{id}\otimes\tau_M)[U(X\otimes I_M)U^*]$ for a unitary $U$ on
$\mathbb C^d\otimes\mathbb C^M$, and $\cK_d$ denotes the closure of the set of channels on $M_d$ with a finite tracial factorization of some bath
dimension. It is also the closure of the convex hull of that set~\cite[Proposition~3.8(ii)]{Musat2025}, and it consists of unital channels. Channels
with a finite tracial factorization are exactly the noisy operations of~\cite{HHO2003} from $M_d$ to $M_d$, and $\cK_d$ is their closure.

\paragraph{Causal services.} A \emph{causal service} of a channel $\Phi$ on $M_d$ for $n$ rounds specifies a \emph{device}, its physical
realization, which serves one use of the channel in each round. The device has a memory register of $Q$ qubits, started in a fixed state
$\beta_0$. In each round $t$ it applies a unitary $U_t$, fixed in advance, to the input and the memory and then releases the $d$-dimensional output. After each round
but the last, a prescribed set of memory qubits is exported for good and replaced by the same number of imported qubits in a fixed state.
That state may be correlated within one step but is independent of everything else. An observer chooses the inputs adaptively, possibly
entangled with references of its own. The \emph{error} of the service is the largest half trace distance between the observer's final state in
this experiment and in the ideal experiment with $n$ independent uses of $\Phi$. Figure~\ref{fig:device} shows the device and its rounds.

The resources are the memory $Q$, the exchange $J$ and the purity $P+C$. The \emph{exchange} is the total number of exported qubits;
exporting is the only way the device can discard a qubit without handing it to the observer. The \emph{purity} is an entropy deficit in bits,
not $\Tr\rho^2$. Here $P=Q-S(\beta_0)$ is the deficit of the initial memory, and $C$ is the sum over the imports of their number of qubits
minus their entropy. A pure qubit is a resource, since it costs work to prepare, while a maximally mixed one is free; the Stinespring device
of a channel with Choi rank $r_*$ spends about $\log r_*$ bits of purity in every round. Every retained qubit counts; there is no free reset,
measurement record or postselection.

% Figure: a causal service (the device) and its rounds.
\begin{figure}[t]
\centering
\begin{tikzpicture}[x=1cm,y=1cm,font=\small,>=Stealth,
  box/.style={draw,fill=white,minimum width=1.1cm,minimum height=1.5cm,inner sep=2pt},
  wire/.style={->,thick},
  xw/.style={->,thick,accent},
  lbl/.style={font=\footnotesize,inner sep=1pt}]
\colorlet{accent}{blue!50!black}
% the memory band
\fill[accent!12] (0,-0.5) rectangle (13.7,0.5);
\draw[accent!60] (0,-0.5) -- (13.7,-0.5) (0,0.5) -- (13.7,0.5);
\node[anchor=west] at (0.2,0) {$\beta_0$};
\node[lbl,anchor=north west,align=left] at (0,-0.65) {memory: $Q$ qubits\\ purity $P=Q-S(\beta_0)$};
% the observer, and the error of the service
\node[draw,rounded corners,align=center,inner ysep=4pt,anchor=south west] (obs) at (2.8,2.62)
  {observer: keeps a quantum reference, chooses each input adaptively\\
   \footnotesize error: half trace distance of its final state from $n$ independent uses of $\Phi$};
% the rounds: the input enters, U_t acts, the output is released before the next input
\foreach \x/\t in {3.5/1,6.7/2,11.2/n}{
  \node[box] at (\x,0.15) {$U_{\t}$};
  \draw[wire] (\x-0.3,2.62) -- (\x-0.3,0.9);
  \draw[wire] (\x+0.3,0.9) -- (\x+0.3,2.62);
  \node[lbl,anchor=east] at (\x-0.42,2.15) {input $\t$};
  \node[lbl,anchor=west] at (\x+0.42,1.35) {output $\t$};
}
\node at (9.05,0.15) {$\cdots$};
\node at (9.05,1.75) {$\cdots$};
% export and import after each round but the last
\foreach \x in {4.7,7.9,9.9}{
  \draw[xw] (\x,-0.5) -- (\x,-1.5);
  \draw[xw] (\x+0.5,-1.5) -- (\x+0.5,-0.5);
}
\node[lbl,anchor=east] at (4.6,-1.0) {export};
\node[lbl,anchor=west] at (5.3,-1.0) {import};
\node[lbl,anchor=north,align=center] at (4.95,-1.65) {exchange $J$: the number of qubits exported, for good\\ purity $C$: qubits minus entropy, summed over the imports};
\end{tikzpicture}
\caption{A causal service of $\Phi$ for $n$ rounds. In each round the input enters, a fixed unitary $U_t$ acts on it and the memory, and the output
is released before the next input arrives; after each round but the last, a prescribed set of memory qubits is exported for good and replaced by
the same number of fresh qubits in a fixed state. The resources are the memory $Q$, the exchange $J$, the purity $P+C$, and the error, the
largest half trace distance between the observer's final state and its final state after $n$ independent uses of $\Phi$.}
\label{fig:device}
\end{figure}


\paragraph{Gadget channels.} Let $G$ be a group generated by a finite set $S$, and let $R$ be a finite set of words in $S^{\pm1}$, each of length at
least two and each representing $1$ in $G$. Write $r_i=\sigma_{i,1}\cdots\sigma_{i,\ell_i}$ ($1\le i\le r_0$) for the words of $R$, and put
$s=|S|$ and $\Lambda=\sum_i\ell_i$. The \emph{symbols} are the $2s$ letters $x^{\pm1}$, $x\in S$, and the proper prefixes
$p_{i,l}=\sigma_{i,1}\cdots\sigma_{i,l}$, $2\le l\le\ell_i-1$, one for each pair $(i,l)$; we also put $p_{i,1}=\sigma_{i,1}$ and let $p_{i,\ell_i}$
be the empty word. The \emph{triangles} are the triples $(p_{i,l-1},\sigma_{i,l},p_{i,l})$, $2\le l\le\ell_i$, and $(x,x^{-1},\varnothing)$,
$x\in S$. For unitaries $\lambda(w)$ assigned to the symbols, with $\lambda(\varnothing)=I$, let
\[
U_\lambda=\mathrm{diag}\bigl(I,(\lambda(w))_{w}\bigr)\ \oplus\ \bigoplus_{(x,y,z)}\frac1{\sqrt2}\begin{pmatrix}I&\lambda(y)\\ \lambda(x)&-\lambda(x)\lambda(y)\end{pmatrix},
\]
where $w$ runs over the symbols and $(x,y,z)$ over the triangles. Each triangle block equals $\mathrm{diag}(I,\lambda(x))\,(H\otimes I)\,\mathrm{diag}(I,\lambda(y))$
with $H$ the Hadamard matrix, so $U_\lambda$ is unitary for all unitary $\lambda$. A triangle is a Hadamard test of one step $xy=z$ of one
relation. Its block is unitary whatever is put in it, and its lower right entry $-\lambda(x)\lambda(y)$ is compared, through the Choi support
defined next, with the entry that carries the label of $z$. A device that reproduces the channel must pass every test with the operators it
holds in memory. There are $2s+\Lambda-2r_0$ symbols and $\Lambda-r_0+s$ triangles, so the system dimension is
\begin{equation}\label{eq:gadget-dimension}
d=1+(2s+\Lambda-2r_0)+2(\Lambda-r_0+s)=1+4s+3\Lambda-4r_0 .
\end{equation}
The \emph{gadget channel} $\Phi_G$ on $M_d$ evaluates each symbol $w$ in the left regular representation, $\lambda(w)=\lambda([w])$ with $[w]\in G$ the
element that $w$ represents, and traces out the group von Neumann algebra $L(G)$ with its canonical trace $\tau$:
$\Phi_G(X)=(\mathrm{id}\otimes\tau)[U_\lambda(X\otimes1)U_\lambda^*]$. In particular $\Phi_G$ is factorizable, through $(L(G),\tau)$. Since
the relations hold, the lower right entry of the block of a triangle
$(x,y,z)$ is $-\lambda([z])$.

We write the gadget as $U_\lambda=\sum_{a\in\mathcal T}c_aE_a\otimes\lambda([w_a])$, where $\mathcal T$ is the set of \emph{slots}, the nonzero
entries, $E_a$ is the matrix unit of slot $a$, and $c_a\in\{1,\pm2^{-1/2}\}$. The word $w_a$ is empty, a symbol, or the concatenation $xy$ in
the lower right corner of a triangle. The elements $[w_a]$ are the \emph{labels}. With $A_g=\sum_{[w_a]=g}c_aE_a$, the matrices $A_g$ have disjoint
supports, $\Phi_G(X)=\sum_gA_gXA_g^*$ and $\mathcal J(\Phi_G)=\frac1d\sum_g|A_g\rangle\!\rangle\langle\!\langle A_g|$, because
$\tau(\lambda(g)\lambda(h)^*)=\delta_{gh}$. The Choi rank $r_*$ of $\Phi_G$ is therefore the number of distinct labels. Every label is carried by an
\emph{anchor}, a diagonal slot with coefficient $1$ whose row and column contain no other slot. The anchor of the label $1$ is the first
diagonal entry, and every other label has a symbol entry as anchor, since the lower right entry of a triangle carries the label of a prefix
symbol or $1$. Feeding one half of a maximally entangled state on the $r_*$ anchor positions gives a pure probe whose output is flat of rank
$r_*$, and $\log r_*$ is the largest entropy of a complementary output of $\Phi_G$~\cite[Proposition~9.1]{DouglasMghirbiB}.

\paragraph{The rank rate.} The probe fixes the least exchange that bounded memory allows. Over $n$ rounds the ideal experiment hands the
observer $n\log r_*$ bits of entropy. The observer's final state is a marginal of a pure state whose other half consists of the final memory,
the exported qubits and the purifiers of the initial memory and of the imports. Those are at most $2(Q+J)$ qubits in all, so a device hands
over at most $2(Q+J)$ bits. With memory sublinear in $n$ it must export about $\frac12\log r_*$ qubits per round at least. This rate is the optimal
exchange rate of a gadget channel in $\cK_d$~\cite[Propositions~3.1 and~9.1]{DouglasMghirbiB}, and we call it the \emph{rank rate}: a service is
at the rank rate if $J\le\frac n2\log r_*$. For the lamp channel, Theorem~\ref{thm:streaming} attains it with memory $n^{1/2+o(1)}$.

\paragraph{One round.} Let $U$ be a unitary on $\mathbb C^d\otimes\mathbb C^M$ and $\rho$ a state on $\mathbb C^M$, and put
\[
\Psi(X)=\mathrm{Tr}_M\,U(X\otimes\rho)U^*,\qquad T(\rho)=\mathrm{Tr}_d\,U(\tau_d\otimes\rho)U^* .
\]
Relative to a channel $\Phi$ with Choi support projection $\Pi_\Phi$, the \emph{support leakage} of $\Psi$ is
$\ell=\Tr[(I-\Pi_\Phi)\mathcal J(\Psi)]$. The \emph{entropy production} of the round is $\Delta S=S(T(\rho))-S(\rho)$, which is nonnegative because
$T$ is unital. A service produces such rounds, and the next result, from~\cite[Theorem~5.2]{DouglasMghirbiB}, finds one in which both quantities are
small. Appendix~\ref{app:imported} reproduces its proof.

\begin{theorem}[Sequential extraction]\label{thm:extraction}
Let a causal service of a channel $\Phi$ on $M_d$ for $n$ rounds have memory $Q$, error $\epsilon<1/2$ and $P+C\le B$, and put
\[
t_n=\frac{B+\epsilon+h_2(\epsilon)+\epsilon\log n}{n(1-\epsilon)} .
\]
Then some round $t$, with its unitary $U_t$ and the memory state $\rho_t$ entering it, defines a channel $\Psi_t(X)=\mathrm{Tr}_M\,U_t(X\otimes\rho_t)U_t^*$
with bath dimension $M=2^Q$, $\ddia(\Psi_t,\Phi)\le\epsilon/(1-\epsilon)$, and support leakage $\ell_t$ and entropy production $\Delta S_t$ satisfying
$\ell_t+\Delta S_t\le t_n$.
\end{theorem}

The bath of this round is the memory itself, of dimension $2^Q$, so $S(\rho_t)\le Q$. Every lower bound on memory in Part~II is a lower bound
on the entropy of a memory state that realizes one good round.

\paragraph{A worked example: the gadget of $S_3$.} Take $G=S_3$ with $S=\{a,b\}$ and $R=\{a^2,\,b^2,\,(ab)^3\}$, so that $s=2$, $r_0=3$ and
$\Lambda=2+2+6=10$. The symbols are the four letters $a^{\pm1},b^{\pm1}$ and the four proper prefixes $ab$, $aba$, $abab$, $ababa$ of
$(ab)^3$; the words $a^2$ and $b^2$ have no proper prefix of length two or more. The triangles are the nine triples $(a,a,\varnothing)$,
$(b,b,\varnothing)$, $(a,b,ab)$, $(ab,a,aba)$, $(aba,b,abab)$, $(abab,a,ababa)$, $(ababa,b,\varnothing)$, $(a,a^{-1},\varnothing)$ and
$(b,b^{-1},\varnothing)$. Formula~\eqref{eq:gadget-dimension} gives $d=1+8+2\cdot9=27$. The labels are the elements of $S_3$ that the slots
carry, namely $1$, $a$, $b$, $ab$, $aba$, $abab=ba$ and $ababa=b$; these are all six elements, so $r_*=6$. Since $S_3$ is finite, its group
von Neumann algebra is the regular representation on $\mathbb C^6$ with the normalized trace. The channel $\Phi_{S_3}$ therefore has a finite
tracial factorization with bath dimension $6$ by definition, and one round with a flat bath on six points reproduces it exactly. Serving $n$
rounds costs little more. The regular action of $S_3$ on itself is a model with no bad points, and Theorem~\ref{thm:streaming} with $k=k'=6$ gives a
service at the rank rate with purity at most $2$ and memory $O(\log n)$. The lamp channel of Section~\ref{sec:lamp-channel} adds five letters
and twenty-six words to these three. What those words cost a device is what Theorem~B measures.
